\section{Supplement on flat modules}
\label{section:0.10}

For the proofs of the properties stated without proof in
n\textsuperscript{os} (\hyperref[subsection:0.10.1]{10.1}) and
(\hyperref[subsection:0.10.2]{10.2}), we refer the reader to Bourbaki,
\emph{Alg. comm.}, chap. II and III.

\subsection{Relations between flat modules and free modules}
\label{subsection:0.10.1}

\begin{env}[10.1.1]
\label{0.10.1.1}
Let $A$ be a ring, $\mathfrak{J}$ an ideal of $A$, and $M$ an $A$-module; for
every integer $p\geq0$, we have a canonical homomorphism of
$(A/\mathfrak{J})$-modules
\[
\label{0.10.1.1.1}
  \varphi_p:(M/\mathfrak{J}M)\otimes_{A/\mathfrak{J}}
  (\mathfrak{J}^p/\mathfrak{J}^{p+1})
  \longrightarrow \mathfrak{J}^pM/\mathfrak{J}^{p+1}M
  \tag{10.1.1.1}
\]
which is evidently \emph{surjective}. We denote by
$\operatorname{gr}(A)=\bigoplus_{p\geq0}\mathfrak{J}^p/\mathfrak{J}^{p+1}$ the graded ring
associated to $A$ filtered by the $\mathfrak{J}^p$, and by
$\operatorname{gr}(M)=\bigoplus_{p\geq0}\mathfrak{J}^pM/\mathfrak{J}^{p+1}M$ the graded
$\operatorname{gr}(A)$-module associated to $M$ filtered by the $\mathfrak{J}^pM$; we
thus have $\operatorname{gr}_p(A)=\mathfrak{J}^p/\mathfrak{J}^{p+1}$,
$\operatorname{gr}_p(M)=\mathfrak{J}^pM/\mathfrak{J}^{p+1}M$; the $\varphi_p$ define
a \emph{surjective} homomorphism of graded $\operatorname{gr}(A)$-modules
\[
\label{0.10.1.1.2}
  \varphi:\operatorname{gr}_0(M)\otimes_{\operatorname{gr}_0(A)}
  \operatorname{gr}(A)\longrightarrow\operatorname{gr}(M).
  \tag{10.1.1.2}
\]
\end{env}

\oldpage[0\textsubscript{III}]{18}
\begin{env}[10.1.2]
\label{0.10.1.2}
Suppose that one of the following hypotheses holds:
\begin{enumerate}
  \item[(i)] $\mathfrak{J}$ is nilpotent;
  \item[(ii)] $A$ is Noetherian, $\mathfrak{J}$ is contained in the radical of
    $A$, and $M$ is of finite type.
\end{enumerate}
Then the following properties are equivalent:
\begin{enumerate}
  \item[$a)$] $M$ is a free $A$-module.
  \item[$b)$] $M/\mathfrak{J}M=M\otimes_A(A/\mathfrak{J})$ is a free
    $(A/\mathfrak{J})$-module, and
    $\operatorname{Tor}_1^A(M,A/\mathfrak{J})=0$.
  \item[$c)$] $M/\mathfrak{J}M$ is a free $(A/\mathfrak{J})$-module and
    the canonical homomorphism (\hyperref[0.10.1.1.2]{10.1.1.2}) is
    injective (hence bijective).
\end{enumerate}
\end{env}

\begin{env}[10.1.3]
\label{0.10.1.3}
Suppose that $A/\mathfrak{J}$ is a \emph{field} (in other words that
$\mathfrak{J}$ is maximal), and that one of the hypotheses (i), (ii) of
(\hyperref[0.10.1.2]{10.1.2}) holds. Then the following properties
are equivalent:
\begin{enumerate}
  \item[$a)$] $M$ is a free $A$-module.
  \item[$b)$] $M$ is a projective $A$-module.
  \item[$c)$] $M$ is a flat $A$-module.
  \item[$d)$] $\operatorname{Tor}_1^A(M,A/\mathfrak{J})=0$.
  \item[$e)$] The canonical homomorphism
    (\hyperref[0.10.1.1.2]{10.1.1.2}) is bijective.
\end{enumerate}
This result will apply in particular in the following two cases:
\begin{enumerate}
  \item[(i)] $M$ is an \emph{arbitrary} module over a local ring $A$ whose
    maximal ideal $\mathfrak{J}$ is \emph{nilpotent} (for example an Artinian
    local ring).
  \item[(ii)] $M$ is a module \emph{of finite type} over a
    \emph{Noetherian} local ring.
\end{enumerate}
\end{env}

\subsection{Local flatness criteria}
\label{subsection:0.10.2}

\begin{env}[10.2.1]
\label{0.10.2.1}
The hypotheses and notation being those of
(\hyperref[0.10.1.1]{10.1.1}), consider the following conditions:
\begin{enumerate}
  \item[$a)$] $M$ is a flat $A$-module.
  \item[$b)$] $M/\mathfrak{J}M$ is a flat $(A/\mathfrak{J})$-module and
    $\operatorname{Tor}_1^A(M,A/\mathfrak{J})=0$.
  \item[$c)$] $M/\mathfrak{J}M$ is a flat $(A/\mathfrak{J})$-module and
    the canonical homomorphism (\hyperref[0.10.1.1.2]{10.1.1.2}) is
    bijective.
  \item[$d)$] For every $n\geq1$, $M/\mathfrak{J}^nM$ is a flat
    $(A/\mathfrak{J}^n)$-module.
\end{enumerate}
We then have the implications
\[
  a\Rightarrow b\Rightarrow c\Rightarrow d
\]
and if $\mathfrak{J}$ is \emph{nilpotent}, the four conditions $a)$, $b)$,
$c)$, $d)$ are \emph{equivalent}. The same holds if $A$ is Noetherian
and moreover $M$ is \emph{ideally separated}, that is, if for every
ideal $\mathfrak{a}$ of $A$, the $A$-module $\mathfrak{a}\otimes_A M$ is
\emph{separated} for the $\mathfrak{J}$-preadic topology.
\end{env}

\begin{env}[10.2.2]
\label{0.10.2.2}
Let $A$ be a Noetherian ring, $B$ a commutative Noetherian $A$-algebra,
$\mathfrak{J}$ an ideal of $A$ such that $\mathfrak{J}B$ is contained in the
radical of $B$, and $M$ a $B$-module of finite type. Then, when $M$ is
considered as an $A$-module, the conditions $a)$, $b)$, $c)$, $d)$ of
(\hyperref[0.10.2.1]{10.2.1}) are equivalent.

This result applies above all when $A$ and $B$ are Noetherian
\emph{local} rings and the homomorphism $A\to B$ is a
\emph{local} homomorphism. More particularly, if $\mathfrak{J}$ is then the
\emph{maximal} ideal of $A$, we can, in conditions $b)$ and $c)$, drop
the hypothesis that $M/\mathfrak{J}M$ is flat, which is automatically satisfied,
and condition $d)$ means that the modules $M/\mathfrak{J}^nM$ are
\emph{free} over the $A/\mathfrak{J}^n$.
\end{env}

\oldpage[0\textsubscript{III}]{19}
\begin{env}[10.2.3]
\label{0.10.2.3}
The hypotheses on $A$, $B$, $\mathfrak{J}$, $M$ being those formulated at the
beginning of (\hyperref[0.10.2.2]{10.2.2}), let $\widehat A$ be the separated
completion of $A$ for the $\mathfrak{J}$-preadic topology, and $\widehat M$ the
separated completion of $M$ for the $\mathfrak{J}B$-preadic topology. Then,
for $M$ to be a flat $A$-module, it is necessary and sufficient that $\widehat M$
be a flat $\widehat A$-module.
\end{env}

\begin{env}[10.2.4]
\label{0.10.2.4}
Let $\rho:A\to B$ be a local homomorphism of Noetherian local rings, $k$
the residue field of $A$, and $M$, $N$ two $B$-modules of finite type, $N$ being
assumed to be $A$-\emph{flat}. Let $u:M\to N$ be a $B$-homomorphism. Then the
following conditions are equivalent:
\begin{enumerate}
  \item[$a)$] $u$ is injective and $\operatorname{Coker}(u)$ is a flat $A$-module.
  \item[$b)$] $u\otimes1:M\otimes_A k\to N\otimes_A k$ is injective.
\end{enumerate}
\end{env}

\begin{env}[10.2.5]
\label{0.10.2.5}
Let $\rho:A\to B$, $\sigma:B\to C$ be local homomorphisms of Noetherian
local rings, $k$ the residue field of $A$, and $M$ a $C$-module of finite
type. Suppose that $B$ is a \emph{flat} $A$-module. Then the following
conditions are equivalent:
\begin{enumerate}
  \item[$a)$] $M$ is a flat $B$-module.
  \item[$b)$] $M$ is a flat $A$-module, and $M\otimes_A k$ is a flat
    $(B\otimes_A k)$-module.
\end{enumerate}
\end{env}

\begin{proposition}[10.2.6]
\label{0.10.2.6}
Let $A$, $B$ be two Noetherian local rings, $\rho:A\to B$ a local
homomorphism, $\mathfrak{J}$ an ideal of $B$ contained in the maximal
ideal, and $M$ a $B$-module of finite type. Suppose that for every $n\geq0$,
$M_n=M/\mathfrak{J}^{n+1}M$ is a flat $A$-module. Then $M$ is a flat
$A$-module.
\end{proposition}

We must prove that for every injective homomorphism $u:N'\to N$ of
$A$-modules of finite type, $v=1\otimes u:M\otimes_A N'\to M\otimes_A N$ is
injective. Now, $M\otimes_A N'$ and $M\otimes_A N$ are $B$-modules of finite
type, hence separated for the $\mathfrak{J}$-preadic topology
(\hyperref[0.7.3.5]{$0_{\mathrm I}$, 7.3.5}); it therefore suffices to prove that
the homomorphism
$\widehat v:\widehat{(M\otimes_A N')}\to\widehat{(M\otimes_A N)}$ of the
separated completions is injective. Now, we have $\widehat v=\varprojlim v_n$, where
$v_n$ is the homomorphism $1\otimes u:M_n\otimes_A N'\to M_n\otimes_A N$;
since by hypothesis $M_n$ is $A$-flat, $v_n$ is injective for every $n$, hence
so is $\widehat v$, the functor $\varprojlim$ being left exact.

\begin{corollary}[10.2.7]
\label{0.10.2.7}
Let $A$ be a Noetherian ring, $B$ a Noetherian local ring,
$\rho:A\to B$ a homomorphism, $f$ an element of the maximal ideal of $B$, and $M$
a $B$-module of finite type. Suppose that the homothety
$f_M:x\mapsto fx$ of $M$ is injective and that $M/fM$ is a flat
$A$-module. Then $M$ is a flat $A$-module.
\end{corollary}

Set $M_i=f^iM$ for $i\geq0$; since $f_M$ is injective,
$M_i/M_{i+1}$ is isomorphic to $M/fM$, hence $A$-flat for every $i\geq0$; from
the exact sequence
\[
  0\longrightarrow M_i/M_{i+1}\longrightarrow M/M_{i+1}
  \longrightarrow M/M_i\longrightarrow0
\]
we deduce by induction on $i$ that $M/M_i$ is $A$-\emph{flat} for every
$i\geq0$ (\hyperref[0.6.1.2]{$0_{\mathrm I}$, 6.1.2}); we can therefore
apply (\hyperref[0.10.2.6]{10.2.6}). We can also argue directly
as follows: for every $A$-module $N$ of finite type, $M\otimes_A N$ is a
$B$-module of finite type; since $f$ belongs to the radical $\mathfrak n$ of
$B$, the $(f)$-adic topology on $M\otimes_A N$ is finer than the
\oldpage[0\textsubscript{III}]{20}
$\mathfrak n$-adic topology, and we know that the latter is \emph{separated}
(\hyperref[0.7.3.5]{$0_{\mathrm I}$, 7.3.5}). Moreover, since $M/M_i$ is
$A$-flat, we have
\[
  f^i(M\otimes_A N)=\operatorname{Im}(M_i\otimes_A N\to M\otimes_A N)
  =\operatorname{Ker}(M\otimes_A N\to(M/M_i)\otimes_A N)
\]
(\hyperref[0.6.1.2]{$0_{\mathrm I}$, 6.1.2}). Now let $N$ be an
$A$-module of finite type, $N'$ a submodule of $N$, and $j:N'\to N$ the canonical
injection; in the commutative diagram
\[
\xymatrix{
  M\otimes_A N'\ar[r]\ar[d]_{1_M\otimes j} &
  (M/M_i)\otimes_A N'\ar[d]^{1_{M/M_i}\otimes j}\\
  M\otimes_A N\ar[r] & (M/M_i)\otimes_A N
}
\]
$1_{M/M_i}\otimes j$ is injective since $M/M_i$ is $A$-flat; we
conclude that
\[
  \operatorname{Ker}(M\otimes_A N'\to M\otimes_A N)
  \subset\operatorname{Ker}(M\otimes_A N'\to(M/M_i)\otimes_A N')
\]
whatever $i$ may be; since the intersection of the right-hand sides reduces to
$0$, as we saw above, the same holds for the left-hand side, and
consequently $M$ is $A$-flat.

\begin{proposition}[10.2.8]
\label{0.10.2.8}
Let $A$ be a reduced Noetherian ring and $M$ an $A$-module of finite type.
Suppose that for every $A$-algebra $B$ which is a discrete valuation
ring, $M\otimes_A B$ is a flat $B$-module (hence free
(\hyperref[0.10.1.3]{10.1.3})). Then $M$ is a flat $A$-module.
\end{proposition}

We know that for $M$ to be flat, it is necessary and sufficient that for every
maximal ideal $\mathfrak m$ of $A$, $M_{\mathfrak m}$ be a flat
$A_{\mathfrak m}$-module
(\hyperref[0.6.3.3]{$0_{\mathrm I}$, 6.3.3}); we may therefore restrict to the case
where $A$ is \emph{local} (\hyperref[0.1.2.8]{$0_{\mathrm I}$, 1.2.8}).
Let then $\mathfrak m$ be the maximal ideal of $A$, $\mathfrak p_i$
$(1\leq i\leq r)$ its minimal prime ideals, and $k$ the residue field
$A/\mathfrak m$. We know (\hyperref[II.7.1.7]{II, 7.1.7}) that for each
$i$ there exists a discrete valuation ring $B_i$ having the same field of
fractions $K_i$ as the integral ring $A/\mathfrak p_i$, and dominating the
latter. Set $M_i=M\otimes_A B_i$. By hypothesis, $M_i$ is free over
$B_i$, so, denoting by $k_i$ the residue field of $B_i$, we have
\[
\label{0.10.2.8.1}
  \operatorname{rg}_{k_i}(M_i\otimes_{B_i}k_i)
  =\operatorname{rg}_{K_i}(M_i\otimes_{B_i}K_i).
  \tag{10.2.8.1}
\]
But it is clear that the composite homomorphism
$A\to A/\mathfrak p_i\to B_i$ is local, so $k_i$ is an extension of $k$,
and we have
$M_i\otimes_{B_i}k_i=M\otimes_A k_i=(M\otimes_A k)\otimes_k k_i$, and on
the other hand $M_i\otimes_{B_i}K_i=M\otimes_A K_i$. Equality
(\hyperref[0.10.2.8.1]{10.2.8.1}) therefore implies
\[
  \operatorname{rg}_k(M\otimes_A k)=\operatorname{rg}_{K_i}(M\otimes_A K_i)
  \qquad\text{for }1\leq i\leq r
\]
and since $A$ is reduced, we know that this condition implies that $M$ is a
\emph{free} $A$-module (Bourbaki, \emph{Alg. comm.}, chap. II, \S~3,
n\textsuperscript{o}~2, prop. 7).

\subsection{Existence of flat extensions of local rings}
\label{subsection:0.10.3}

\begin{proposition}[10.3.1]
\label{0.10.3.1}
Let $A$ be a Noetherian local ring, $\mathfrak J$ its maximal ideal, and
$k=A/\mathfrak J$ its residue field. Let $K$ be an extension of the field $k$. There
exists a local homomorphism from $A$ to a Noetherian local ring $B$ such
that $B/\mathfrak J B$ is $k$-isomorphic to $K$ and $B$ is a flat
$A$-module.
\end{proposition}

We will prove this proposition in several steps.

\begin{env}[10.3.1.1]
\label{0.10.3.1.1}
Suppose first that $K=k(T)$, where $T$ is an indeterminate. In the
polynomial ring $A'=A[T]$, consider the prime ideal
$\mathfrak J'=\mathfrak J A'$ consisting of the
\oldpage[0\textsubscript{III}]{21}
polynomials having their coefficients in the ideal $\mathfrak J$; it is clear
that $A'/\mathfrak J'$ is canonically isomorphic to $k[T]$. Let us show that
the ring of fractions $B=A'_{\mathfrak J'}$ answers the question; it is
evidently a Noetherian local ring with maximal ideal
$\mathfrak L=\mathfrak J B$. Moreover,
$B/\mathfrak L=(A'/\mathfrak J')_{\mathfrak J'}=(k[T])_{\mathfrak J'}$
is none other than the field of fractions $K$ of $k[T]$. Finally, $B$ is a
flat $A'$-module and $A'$ a free $A$-module, hence $B$ is a flat $A$-module
(\hyperref[0.6.2.1]{$0_{\mathrm I}$, 6.2.1}).
\end{env}

\begin{env}[10.3.1.2]
\label{0.10.3.1.2}
Suppose next that $K=k(t)=k[t]$, where $t$ is algebraic over $k$; let
$f\in k[T]$ be the minimal polynomial of $t$; there exists a monic polynomial
$F\in A[T]$ whose canonical image in $k[T]$ is $f$. Again set
$A'=A[T]$, and let $\mathfrak J'$ be the ideal $\mathfrak J A'+(F)$ in $A'$.
We will see that the quotient ring $B=A'/(F)$ answers the question this
time. First of all, it is clear that $B$ is a \emph{free} $A$-module,
hence flat. The ring $A'/\mathfrak J'$ is isomorphic to
\[
  (A'/\mathfrak J A')/((\mathfrak J A'+(F))/\mathfrak J A')
  =k[T]/(f)=K~;
\]
the image $\mathfrak L$ of $\mathfrak J'$ in $B$ is therefore maximal, and we
evidently have $\mathfrak L=\mathfrak J B$. Finally, $B$ is a semi-local ring,
since it is an $A$-module of finite type (Bourbaki, \emph{Alg. comm.}, chap.
IV, \S~2, n\textsuperscript{o}~5, cor. 3 of prop. 9), and its maximal ideals correspond bijectively to those of $B/\mathfrak J B$
([13], vol. I, p. 259); what precedes therefore proves that $B$ is a
local ring.
\end{env}

\begin{lemma}[10.3.1.3]
\label{0.10.3.1.3}
Let $(A_\lambda,f_{\mu\lambda})$ be a filtered inductive system of local
rings such that the $f_{\mu\lambda}$ are local homomorphisms; let
$\mathfrak m_\lambda$ be the maximal ideal of $A_\lambda$, and let
$K_\lambda=A_\lambda/\mathfrak m_\lambda$. Then
$A'=\varinjlim A_\lambda$ is a local ring of which
$\mathfrak m'=\varinjlim\mathfrak m_\lambda$ is the maximal ideal and
$K=\varinjlim K_\lambda$ the residue field. Moreover, if
$\mathfrak m_\mu=\mathfrak m_\lambda A_\mu$ for $\lambda<\mu$, we have
$\mathfrak m'=\mathfrak m_\lambda A'$ for every $\lambda$. If, in addition, for
$\lambda<\mu$, $A_\mu$ is a flat $A_\lambda$-module, and if all the
$A_\lambda$ are Noetherian, then $A'$ is Noetherian and is a flat
$A_\lambda$-module for every $\lambda$.
\end{lemma}

Since $f_{\mu\lambda}(\mathfrak m_\lambda)\subset\mathfrak m_\mu$ for
$\lambda<\mu$ by hypothesis, the $\mathfrak m_\lambda$ form an inductive
system, and its limit $\mathfrak m'$ is evidently an ideal of $A'$.
Moreover, if $x'\notin\mathfrak m'$, there exists $\lambda$ such that
$x'=f_\lambda(x_\lambda)$ for some $x_\lambda\in A_\lambda$
($f_\lambda:A_\lambda\to A'$ denoting the canonical homomorphism); since
$x'\notin\mathfrak m'$, we necessarily have
$x_\lambda\notin\mathfrak m_\lambda$, so $x_\lambda$ has an inverse
$y_\lambda$ in $A_\lambda$, and $y'=f_\lambda(y_\lambda)$ is the inverse of
$x'$ in $A'$, which proves that $A'$ is a local ring with maximal ideal
$\mathfrak m'$; the assertion concerning $K$ follows immediately from the fact
that $\varinjlim$ is an exact functor. The hypothesis
$\mathfrak m_\mu=\mathfrak m_\lambda A_\mu$ means that the canonical
map
$\mathfrak m_\lambda\otimes_{A_\lambda}A_\mu\to\mathfrak m_\mu$ is
surjective; the relation $\mathfrak m'=\mathfrak m_\lambda A'$ therefore again
follows from the exactness of the functor $\varinjlim$ and from the fact that it
commutes with the tensor product.

Suppose now that for $\lambda<\mu$ we have
$\mathfrak m_\mu=\mathfrak m_\lambda A_\mu$ and that $A_\mu$ is a flat
$A_\lambda$-module. Then $A'$ is a flat $A_\lambda$-module for every
$\lambda$, by virtue of (\hyperref[0.6.2.3]{$0_{\mathrm I}$, 6.2.3});
since $A'$ and $A_\lambda$ are local rings and
$\mathfrak m'=\mathfrak m_\lambda A'$, $A'$ is even a
\emph{faithfully flat} $A_\lambda$-module
(\hyperref[0.6.6.2]{$0_{\mathrm I}$, 6.6.2}). Finally, suppose in addition
that the $A_\lambda$ are \emph{Noetherian}; the
$\mathfrak m_\lambda$-preadic topologies are then separated
(\hyperref[0.7.3.5]{$0_{\mathrm I}$, 7.3.5}); let us show that it follows
first of all that on $A'$ the $\mathfrak m'$-adic topology is \emph{separated}.
Indeed, if $x'\in A'$ belongs to all the $\mathfrak m'^{n}$ $(n>0)$,
it is the image of some $x_\mu\in A_\mu$ for a certain index $\mu$, and since
the inverse image in $A_\mu$ of
$\mathfrak m'^{n}=\mathfrak m_\mu^n A'$ is $\mathfrak m_\mu^n$
(\hyperref[0.6.6.1]{$0_{\mathrm I}$, 6.6.1}), $x_\mu$ belongs to all
the $\mathfrak m_\mu^n$, hence $x_\mu=0$ by hypothesis, and consequently
$x'=0$. Let us denote by $\widehat A'$ the completion of $A'$ for the
$\mathfrak m'$-adic topology; what precedes shows that we have
$A'\subset\widehat A'$. We will show that $\widehat A'$ is
\emph{Noetherian} and $A_\lambda$-\emph{flat} for every $\lambda$; it will
\oldpage[0\textsubscript{III}]{22}
follow that $\widehat A'$ is $A'$-flat
(\hyperref[0.6.2.3]{$0_{\mathrm I}$, 6.2.3}), and since
$\mathfrak m'\widehat A'\ne\widehat A'$, $\widehat A'$ is a faithfully
flat $A'$-module (\hyperref[0.6.6.2]{$0_{\mathrm I}$, 6.6.2}), whence we
will finally conclude that $A'$ is \emph{Noetherian}
(\hyperref[0.6.5.2]{$0_{\mathrm I}$, 6.5.2}), which will complete the
proof of the lemma.

We have $\widehat A'=\varprojlim_n A'/\mathfrak m'^{n}$; because
$A'$ is $A_\lambda$-flat, we have
\[
  \mathfrak m'^{n}/\mathfrak m'^{n+1}
  =(\mathfrak m_\lambda^n/\mathfrak m_\lambda^{n+1})
    \otimes_{A_\lambda}A'
  =(\mathfrak m_\lambda^n/\mathfrak m_\lambda^{n+1})
    \otimes_{K_\lambda}(K_\lambda\otimes_{A_\lambda}A')
  =(\mathfrak m_\lambda^n/\mathfrak m_\lambda^{n+1})
    \otimes_{K_\lambda}K~;
\]
since $\mathfrak m_\lambda^n/\mathfrak m_\lambda^{n+1}$ is a
$K_\lambda$-vector space of finite dimension,
$\mathfrak m'^{n}/\mathfrak m'^{n+1}$ is a $K$-vector space of finite
dimension for every $n\geq0$. It therefore follows from
(\hyperref[0.7.2.12]{$0_{\mathrm I}$, 7.2.12}) and
(\hyperref[0.7.2.8]{$0_{\mathrm I}$, 7.2.8}) that $\widehat A'$ is
Noetherian. We know moreover that the maximal ideal of $\widehat A'$ is
$\mathfrak m'\widehat A'$ and that
$\widehat A'/\mathfrak m'^{n}\widehat A'$ is isomorphic to
$A'/\mathfrak m'^{n}$; since
$A'/\mathfrak m'^{n}=(A_\lambda/\mathfrak m_\lambda^n)
\otimes_{A_\lambda}A'$, $A'/\mathfrak m'^{n}$ is a flat
$(A_\lambda/\mathfrak m_\lambda^n)$-module
(\hyperref[0.6.2.1]{$0_{\mathrm I}$, 6.2.1}); the criterion
(\hyperref[0.10.2.2]{10.2.2}) is therefore applicable to the Noetherian
$A_\lambda$-algebra $\widehat A'$, and shows that $\widehat A'$
is $A_\lambda$-flat. Q.E.D.

\begin{env}[10.3.1.4]
\label{0.10.3.1.4}
Let us now tackle the general case. There exist an ordinal $\gamma$ and, for
every ordinal $\lambda\leq\gamma$, a subfield $k_\lambda$ of $K$ containing
$k$, such that: $1^\circ$ For every $\lambda<\gamma$, $k_{\lambda+1}$ is
an extension of $k_\lambda$ generated by a single element; $2^\circ$ For
every ordinal $\mu$ without a predecessor,
$k_\mu=\bigcup_{\lambda<\mu}k_\lambda$; $3^\circ$ $K=k_\gamma$. Indeed, it
suffices to consider a bijection $\xi\mapsto t_\xi$ of the set of
ordinals $\xi\leq\beta$ (for a suitable $\beta$) onto $K$, to define
$k_\lambda$ by transfinite induction (for $\lambda\leq\beta$) as the
union of the $k_\mu$ for $\mu<\lambda$ if $\lambda$ has no predecessor,
and, if $\lambda=\nu+1$, as $k_\nu(t_\xi)$, where $\xi$ is the smallest
ordinal such that $t_\xi\notin k_\nu$; $\gamma$ is then by definition the
smallest ordinal $\leq\beta$ such that $k_\gamma=K$.

That being so, we will define, by transfinite induction, a family
of Noetherian local rings $A_\lambda$ for $\lambda\leq\gamma$, and local
homomorphisms $f_{\mu\lambda}:A_\lambda\to A_\mu$ for
$\lambda\leq\mu$, satisfying the following conditions:
\begin{enumerate}
  \item[(i)] $(A_\lambda,f_{\mu\lambda})$ is an inductive system and $A_0=A$.
  \item[(ii)] For every $\lambda$, we have a $k$-isomorphism
    $A_\lambda/\mathfrak J A_\lambda\xrightarrow{\sim}k_\lambda$.
  \item[(iii)] For $\lambda\leq\mu$, $A_\mu$ is a flat $A_\lambda$-module.
\end{enumerate}

Suppose then that the $A_\lambda$ and the $f_{\mu\lambda}$ have been defined for
$\lambda<\mu<\xi$, and suppose first that $\xi=\zeta+1$, so
that $k_\xi=k_\zeta(t)$. If $t$ is transcendental over $k_\zeta$, we define
$A_\xi$, following the procedure of (\hyperref[0.10.3.1.1]{10.3.1.1}), to be
equal to $(A_\zeta[t])_{\mathfrak J A_\zeta[t]}$;
$f_{\xi\zeta}$ is the canonical map, and for $\lambda<\zeta$ we take
$f_{\xi\lambda}=f_{\xi\zeta}\circ f_{\zeta\lambda}$; the verification of
conditions (i) to (iii) is then immediate, in view of what was proved in
(\hyperref[0.10.3.1.1]{10.3.1.1}). Next suppose $t$ algebraic, and
let $h$ be its minimal polynomial in $k_\zeta[T]$ and $H$ a monic polynomial in
$A_\zeta[T]$ whose image in $k_\zeta[T]$ is $h$; we then take $A_\xi$
equal to $A_\zeta[T]/(H)$, the $f_{\xi\lambda}$ being defined as
before; the verification of conditions (i) to (iii) then follows from
what was seen in (\hyperref[0.10.3.1.2]{10.3.1.2}).

Now suppose that $\xi$ has no predecessor; we then take
for $A_\xi$ the inductive limit of the inductive system of local rings
$(A_\lambda,f_{\mu\lambda})$ for $\lambda<\xi$; $f_{\xi\lambda}$ is defined
as the canonical map for $\lambda<\xi$. The fact that $A_\xi$ is
local and Noetherian, that the $f_{\xi\lambda}$ are local homomorphisms,
and conditions (i) to (iii) for $\lambda\leq\xi$
\oldpage[0\textsubscript{III}]{23}
then follow from the induction hypothesis and from Lemma
(\hyperref[0.10.3.1.3]{10.3.1.3}). Once this construction is made, it is clear
that the ring $B=A_\gamma$ satisfies the statement of
(\hyperref[0.10.3.1]{10.3.1}).
\end{env}

Note that, by virtue of (\hyperref[0.10.2.1]{10.2.1, $c)$}), we have a
canonical isomorphism
\[
\label{0.10.3.1.5}
  \operatorname{gr}(A)\otimes_k K\xrightarrow{\sim}\operatorname{gr}(B).
  \tag{10.3.1.5}
\]

On the other hand, we can replace $B$ by its $\mathfrak J B$-adic completion
$\widehat B$ without changing the conclusions of
(\hyperref[0.10.3.1]{10.3.1}), since $\widehat B$ is a flat $B$-module
(\hyperref[0.7.3.3]{$0_{\mathrm I}$, 7.3.3}), hence a flat $A$-module
(\hyperref[0.6.2.1]{$0_{\mathrm I}$, 6.2.1}).

We have moreover proved the

\begin{corollary}[10.3.2]
\label{0.10.3.2}
If $K$ is an extension of finite degree, we may assume that $B$ is a
finite $A$-algebra.
\end{corollary}

